The classical reputation hypothesis predicts repair after social stigma.
Agents who are named or scored poorly should give more to restore standing \pcite{milinski2002reputation,wedekind2000cooperation}.
Our multi-seed evidence rejects that prediction for developing agents assigned to the non-enforcing institution.
Across both independent realizations, mean next-round intensity falls after gossip and bad-reputation events in SFI.
Coded rationales in Seed~1 show opportunistic framing far more often than reputation management.
Social monitoring therefore associates with withdrawal where Stage-2 enforcement is unavailable.

SI responses reverse sign across seeds.
Enforcing-institution agents can raise intensity after gossip in one realization and lower it in another.
That heterogeneity is scientifically informative.
Social-information effects are contingent on institutional role and stochastic path.
Claims that reputation always stabilizes LLM cooperation overclaim relative to this protocol.
Multi-seed reporting is what makes the qualification possible.
A single negative SI gossip mean would have invited overgeneralization.

Positive average cooperation leaves the repair hypothesis unsupported.
Both seeds sustain moderately positive mean intensity without strong path stability or fairness-heavy contribution talk.
Seed~2's higher mean intensity even coincides with a larger terminal wealth gap.
The pattern matches warnings that LLM reasoning can sustain calculated free-riding alongside occasional giving \pcite{li2025spontaneous,freeriders2025}.
It also aligns with limited-norm findings in multi-agent LLM studies \pcite{gupta2025norms,ashery2025conventions}.
Transfers are real.
Internalized contribution norms are not established by average levels alone.

Forced assignment is essential to interpretation.
Because developed agents cannot exit SI and developing agents cannot enter Stage-2, social marks cannot be answered by institutional flight of the kind observed in free-choice baselines.
Withdrawal becomes a residual response for SFI agents who are named but cannot punish.
Identification without sanction reach, already visible in mixed baselines, reappears here as stigma without repair.
Ambient low consistency scores make gossip eligibility common, yet repair language remains rare.
Naming is therefore cheap relative to behavioral correction.

Democratic adaptation adds a second layer of contingency.
Both seeds expand subsidies and equity weights, yet only Seed~2 weakens punishment.
Governance is endogenous and path dependent under identical whitelists.
That divergence cautions against treating any one political trajectory as the climate society's institutional destiny.
It also confounds shock attribution, because votes and damages share a calendar.
Costless meta-governance can substitute for costly peer enforcement without guaranteeing cooperative norms.

Fiscal outcomes reinforce the political-economy punchline.
Identical coverage ratios coexist with different terminal pools and larger wealth gaps in the higher-cooperation seed.
Collection is not equalization.
A hidden dual-use pool can grow while stock inequality widens.
Mechanism laboratories of this kind therefore speak to design tensions in climate risk sharing without evaluating any real-world fund's effectiveness.
The practical implication for LLM multi-agent design is narrower and sharper.
Social information modules should be evaluated with event studies and multi-seed sign checks, especially when agents lack exit and Stage-2 reach.
